\cotefile{157-5}
\scanpage{437}
\mspage{349}
% slanted
\margin{À rejeter plus bas, après 4.13.31, qui \unsure{est|donne} la première \unsure{application} importante de 4.13.22, alors que 4.13.28 se prouve sans avoir à utiliser 4.13.22 \unsure{ou|ni} ses variantes.}
\uline{Proposition} 4.13.28. \unsure{Soient|\deleted{Soient}} $M,\pi,\Phi$ comme dans \deleted{\unsure{4}} 4.13.23. Supposons de plus que \deleted{\unsure{$\Phi$}} $J\in\Phi$ implique que $J$ est \uline{filtrante}. Alors l'ens $\widetilde{\Sigma}_{\pi}^{\Phi}$ \unsure{des|\deleted{des}} sous-ensembles $\Phi$-$\pi$-accessibles de $M$ est stable par \deleted{\unsure{réunions}} réunions finies. \deleted{Supposons} \inserted{Plus \unsure{gén$^{\rm t}$},} soit $\lambda$ un cardinal \inserted{$<\pi$} tel que \deleted{\unsure{$\Phi$}} $\forall\,J\in\Phi$, $J$ soit grde devant $\lambda$. Alors l'ens $\widetilde{\Sigma}_{\pi}^{\Phi}$ est stable par réunions \deleted{\unsure{index}} $(S_{\unsure{r|i}})_{r\in A}$ \unsure{indexées|indiciées} par des ens. $A$ de cardinal $\leq\lambda$.

\uline{Corollaire} 4.13.29. L'ens des parties accessibles de $\mathrm{Ob}\,M$ est stable par \inserted{``petites''} réunions --- i.e. si $(S_r)_{r\in A}$ est une \inserted{\unsure{petite}} famille de parties accessibles de $\mathrm{Ob}\,M$, alors $\bigcup_{r\in A}S_r$ est aussi accessible.

La démonstration de 4.13.28 repose sur le lemme suivant (je n'explicite pas les démonstrations):
% archive sheet number: 436
\scanpage{438}
\mspage{350}
\uline{Lemme} 4.13.30. Soit $I$ une catégorie grde devant $\pi$, et $(J_{\unsure{r|i}})_{r\in A}$ une famille de sous-catégories pleines telles que $\mathrm{Card}\,A\leq\pi$, \deleted{et que $\mathrm{Ob}\,I=\bigcup_{r\in A}\mathrm{Ob}\,J_r$. (ou plus \unsure{généralt|généralement},} telles que $\bigcup_{r\in A}\mathrm{Ob}\,J_r$ soit cofinal dans $\mathrm{Ord}(I)$, i.e pour tt $i\in\mathrm{Ob}\,I$ \deleted{\unsure{et}} $\exists\,j\in\bigcup_{r\in A}\mathrm{Ob}\,J_r$, et une flèche $i\unsure{\text{---}|\longrightarrow}j$). Alors il existe \deleted{un} $r\in A$ tel\deleted{\unsure{le}} que $J_r$ soit cofinale dans $I$.

\deleted{\uline{\unsure{Proposition}}} \inserted{Théorème} 4.13.31. Soient $M,N$ des catégories accessibles, $f:M\unsure{\text{---}|\longrightarrow}N$ un foncteur accessible, \deleted{\unsure{$S$,}} $T\subset\mathrm{Ob}\,N$ une partie \deleted{accessible} de $N$, $\underline T\subset N$ la sous-catégorie str. pleine \unsure{correspondante} de $N$, \deleted{accessible dans $N$. Alors} $S=f^{-1}(T)$, $\underline S=f^{-1}(\underline T)$ la sous-catégorie pleine de $M$ associée. \unsure{Si} $T$ est accessible, alors $S$ est accessible\,; en d'autres termes, si $\underline T$ est accessible \unsure{en|dans} $N$, $\underline S$ est accessible dans $M$.

Cel\unsure{a} résulte de l'énoncé plus précis\unsure{:}

\uline{Corollaire} 4.13.32. Soient $\pi_0$ \inserted{et $\pi\unsure{>|\geq}\pi_0$} deux cardinaux ayant les propriétés suivantes :

\deleted{\unsure{1°} Il existe un \deleted{cardinal} ordinal limite $\alpha_0$, avec $\mathrm{Card}\,\alpha_0\leq\pi$, et tel que $\underline T$ soit $\alpha_0$-$\pi$-accessible dans $N$.}
% archive sheet number: 437
\scanpage{439}
\mspage{351}
a) $\underline T$ est $\pi_0$-$\pi$-accessible dans $N$ (4.13.23)

b) $\pi_0$ est utile pour $M$, \inserted{(il l'est pour $N$, \unsure{à cause de a)})} \deleted{et \unsure{pour $N$|$f$ est} $\pi_0$-accessible.}

% slanted
\inserted{c) $f(M_\pi)\subset N_\pi$}
% slanted
\margin{\deleted{d) $f(M_{\unsure{\pi|\pi_0}})\unsure{\subset|\in}N_{\unsure{\pi|\pi_0}}$}}
d) \deleted{$f$ est $\pi_0$-\unsure{accessible}.} $M$ satisfait à $L_{\pi_0}$, et $f$ est $\pi_0$-accessible.

\deleted{Alors d) $f$ est $\pi_0$-accessible.}

\deleted{\unsure{(}On peut choisir \unsure{pour} $\pi_0$ un cardinal qui soit utile pour $M$ et pour $N$, et tel que $f$ soit $\pi_0$-} \deleted{Alors} (L'existence d'un tel couple $(\pi_0,\pi)$ résulte aussitôt de 4.13.16. \inserted{et de \unsure{4.12.1|4.22.1}}) Alors $\underline S$ est $\pi_0$-$\pi$-accessible dans $N$.

Comme $f$ commute aux $\underrightarrow{\lim}$ grdes devant $\pi_0$ en vertu de d), et que $\underline T$ est stable par ces limites par a), (cf 4.13.23, 3°), il en est donc de $=$ de $\underline S$, d'où le 3° de la déf. 4.13.23 est satisfait. D'autre part, \deleted{\unsure{comme}} la condition 1° ( \deleted{\unsure{et}} $\pi$ utile pour $M$, $M$ stable par $\underrightarrow{\lim}$ grdes devant $\pi_0$ et de cardinal $\leq\pi$) est vérifiée par b) et d). Donc il reste à vérifier la condition 2°), i.e. que $\underline S$ est $\pi$-accessible. \deleted{on voit que} \deleted{En vertu de 4.13.18} \deleted{Posons} Posons
\[
\begin{aligned}
\underline T(\pi)&=\underline T\cap N_\pi\\
\underline S(\pi)&=\underline S\cap M_\pi\underset{\text{grâce à c)}}{\subset}f^{-1}(\underline T(\pi))
\end{aligned}
\]
\deleted{\unsure{A}} Comme $\underline S$ est stable \inserted{dans $M$} par $\underrightarrow{\lim}$ grdes devant $\pi_0$, \inserted{lesquelles existent dans $M$ par d)} et \unsure{a fortiori|à fortiori} par $\underrightarrow{\lim}$ grdes devant $\pi$, \deleted{\unsure{les}} on a un foncteur canonique pl. fid.
\[
\mathrm{Ind}_{\pi}\,\underline S(\pi)\longrightarrow\underline S,
\]
% archive sheet number: 438
\scanpage{440}
\mspage{352}
D'ailleurs, $\underline T$ étant str. pleine dans $N$, $\underline S$ est str. pleine dans $M$, et il reste à vérifier que le foncteur précédent est ess. surjectif. Soit donc $x\in\mathrm{Ob}\,\underline S$. On aura, grâce à b)
\[
x=\underset{I}{\underrightarrow{\lim}}\,x_i
\qquad\left|\begin{aligned}
 &x_i\in\deleted{\unsure{S}}M_{\deleted{\unsure{<}}\pi}\\
 &I=M_\pi/x
\end{aligned}\right. ,
\]
\deleted{et \unsure{de} foncteur $i\longmapsto x_i$, $I\longrightarrow M_\pi$}

D'ailleurs, le système inductif précédent est $\pi_0$-$\pi$-adapté, i.e (p. 339)

1) $I$ équiv. à une petite catégorie, $I$ grde devant $\pi$ (résulte du fait que $\pi$ est utile pour $M$, cf \textcircled{b}.

2) les limites inductives grdes devant $\pi_0$ et de cardinal $\leq\unsure{\pi|2}$ existent dans $I$ \inserted{résulte de $L_{\pi_0}$ sur $M$ [cf. \textcircled{d}]}

3) le foncteur $\xi:I\unsure{\text{---}|\longrightarrow}M$ y commute.

4) Les $x_i$ sont dans $M_\pi$.

Comme $f$ commute aux $\underrightarrow{\lim}$ grdes devant $\pi_0$, \deleted{\unsure{et}} (cf \textcircled{d}), a fortiori à celles grdes devant $\pi$, on \unsure{a|aura} donc
\[
f(x)\simeq\underset I{\underrightarrow{\lim}}\,f(x_i),
\]
je dis que le syst. ind. $(f(x_i))_{i\in I}$ de $N$ est $\pi_0$-$\pi$-adapté. Les conditions 1), 2) \unsure{sont en évidence} (ne concernant que $I$), le foncteur $I\longrightarrow N$ est le composé
% slanted
\margin{Cette commutation devrait avoir un nom, \unsure{disons|la} condition $L(\pi_0,\pi)$ sur un foncteur}
\diagram{\begin{tikzcd}[ampersand replacement=\&,column sep=small] I\arrow[r,"\xi"]\&M\&\overset{f}{\unsure{\text{---}|\longrightarrow}}\&N \end{tikzcd}}, donc commute aux $\underrightarrow{\lim}$ grdes devant $\pi_0$ et de cardinal $\leq\pi$, puisqu'il en est ainsi de $\xi$ et de $f$, enfin les $f(x_i)$ sont ds $N_\pi$ puisque par \textcircled{c} $f(M_{\unsure{\pi|\pi_0}})\subset N_\pi$.
% archive sheet number: 439
\scanpage{441}
\mspage{353}
Par ailleurs $f(x)\in\mathrm{Ob}\,\underline T$ (puisque $x\in\underline{\mathrm{Ob}}\,\underline S$, \unsure{$\therefore$|et} $\underline S\overset{\mathrm{def}}{=}f^{-1}(\underline T)$). Donc on peut appliquer 4.13.22 $a\Longrightarrow c$ à la situation $(N,\pi_0,\pi,\underline S)$ et au système inductif $(f(x_i))_{i\in I}$, pour conclure que la sous-catégorie $I'$ de $I$, formée des $i$ tels que $f(x_i)\in\deleted{\unsure{\mathrm{Ob}}}\underline S\deleted{\unsure{S}}$ i.e \unsure{alors|tels} que $x_i\in\underline S$ donc $x_i\in\underline S(\pi)$, est \uline{cofinale} dans $I$. Donc on a
\[
x\simeq\underset{i\in I'}{\underrightarrow{\lim}}\,x_{i\unsure{\deleted{{}'}|{}'|\ast}},
\]
avec cette fois les $x_{i\deleted{{}'}}$ dans $\underline S(\pi)$, et $I'$ grde devant $\pi$ (car cofinale \inserted{et pleine} dans $I$, qui est grde devant $\pi$), ce qui prouve que $x$ est ds l'image essentielle de $\mathrm{Ind}_\pi(\underline S(\pi))$, qed.

\uline{Questions} 4.13.33. On vient de voir que si $f:M\unsure{\text{---}|\longrightarrow}N$ est un foncteur accessible entre catégories accessibles, alors l'image inverse d'une \deleted{\unsure{sous-catég}} partie accessible de $\mathrm{Ob}\,N$ est une partie accessible de $\mathrm{Ob}\,M$.
% slanted
\margin{Faux, $=$ si on prend la \unsure{$G$-clôture|$\mathcal G$-clôture} de l'image essentielle (``image $G$-essentielle'')}
Je conjecture que l'énoncé \unsure{similaire|symétrique}, \inserted{non pas} pour les images directes, mais pour les \uline{images essentielles}, est
% archive sheet number: 440
\scanpage{442}
\mspage{354}
vraie, i.e que pour toute \deleted{partie} sous-catégorie $\underline S$ de $M$ accessible dans $M$, \deleted{\unsure{$f$}} son image essentielle dans $N$ est accessible dans $N$. Dans cette question, $\mathrm{OPS}$ que $\underline S=M$, et la question revient à celle-ci : tt foncteur accessible \inserted{$f:M\unsure{\text{---}|\longrightarrow}N$} entre catégories accessibles se factorise-t-il en un foncteur (accessible) \uline{ess. surjectif}, et un foncteur accessible \uline{pl. fidèle} \diagram{\begin{tikzcd}[ampersand replacement=\&,column sep=small] M\arrow[r]\&\underline T\arrow[r]\&N \end{tikzcd}} (où $\underline T$ est elle-$=$ accessible) ?
% slanted
\margin{Il faut se borner tout au moins à ess. \unsure{$G$-surjectif|$\mathcal G$-surjectif}, mais $=$ cel\unsure{à|a} est très optimiste.}
Si oui, cela me paraîtrait peut-être le résultat le plus profond (peut-être le seul \unsure{véritablemt|véritablement} profond) en théorie des catégories.

Voici une question qui me paraît très voisine. Soient $M$ une catégorie accessible, $A$ une sous-catégorie pleine petite, Supposons que $M$ satisfasse à $L_{\pi_0}$ (p. ex. $\pi_0$ est utile pour $M$), et soit $\underline S$ la plus petite sous-catégorie str. pleine de $M$ contenant $A$, et stable par limites grdes devant $\pi_0$. \deleted{Alors} Je conjecture que $\underline S$ est accessible dans $M$.
% archive sheet number: 441
\scanpage{443}
\mspage{355}
De $=$, si $\Phi$ est un \deleted{\illegible{1}}\inserted{petit} ens. de \inserted{petites} catégories, \deleted{et si $\Phi$ est \unsure{assez} petit \unsure{en}} \deleted{\unsure{le|et}} tel que $M$ soit stable par $\underrightarrow{\lim}$ de type $\Phi$, on peut prendre l'enveloppe \inserted{$\underline S$} de $A$ pour les $\underrightarrow{\lim}$ grdes devant $\pi_0$ \uline{et} les $\underrightarrow{\lim}$ de type $\Phi$, et je conjecture encore que $\underline S$ est accessible. Dans cette direction, on a du moins le résultat très partiel suivant :

\uline{Proposition} 4.13.34. Soit $M$ une catégorie accessible stable par \inserted{petites} $\underrightarrow{\lim}$ filtrantes, et soit $\underline S$ une sous-catégorie str. pleine de $M$. Conditions équivalentes :

a) $\underline S$ est accessible dans $M$, et stable dans $M$ par \inserted{petites} $\underrightarrow{\lim}$ filtrantes.

b) Il existe une petite \deleted{\unsure{partie}} \inserted{\unsure{sous-catég}} str. pleine $A$ de $M$, telle que $\underline S$ soit la plus petite sous-catégorie str. pleine de $M$ \unsure{contenant|cont$^{\rm nt}$} $A$, et stable par petites $\underrightarrow{\lim}$ filtrantes.
% archive sheet number: 442
\scanpage{444}
\mspage{356}
\unsure{$\cdot$|.} On a a) $\Longrightarrow$ b), \deleted{\unsure{plus}} car si $\underline S$ est $\pi$-accessible dans $M$, il est immédiat (sans autre condition de stabilité sur $\underline S$) que $\underline S$ est la plus petite sous-catégorie str. pleine de $M$ contenant $\underline S(\pi)\overset{\mathrm{def}}{=}\underline S\cap M_\pi$, et stable par $\underrightarrow{\lim}$ grdes devant $\pi$. Si donc on sait déjà que $\underline S$ est $=$ stable par \uline{ttes} $\underrightarrow{\lim}$ filtrantes petites, il s'ensuit que $\underline S$ est la plus petite sous catégorie de $M$ ayant cette propriété de stabilité plus forte, et \unsure{contenant|cont$^{\rm nt}$} $A$.

Prouvons $b\Longrightarrow a$. Soit $\pi$ un cardinal utile pour $M$, et tel que $A\subset M_\pi$. Soit $B$ la plus petite sous-catégorie str. pleine de $M_{\unsure{\pi|\pi_0}}$, \deleted{stable,} contenant $A$ et stable par $\underrightarrow{\lim}$ filtrantes de cardinal $\leq\pi$. On a donc $B\subset M_\pi$ et $B\subset S$, \deleted{\unsure{donc}} et $\underline S$ est aussi la $L$-enveloppe de $B$. D'ailleurs $B$ est stable par
% archive sheet number: 443
\scanpage{445}
\mspage{357}
facteurs directs. Considérons donc l'image essentielle \inserted{$\underline S'$} de $\mathrm{Ind}_\pi(B)$ dans $M$. On a $\underline S'\subset\underline S$, et \unsure{tout} revient à prouver que $\underline S'=\underline S$, \deleted{cel\unsure{a} prouve} \inserted{d'où} résultera en effet que $\underline S$ est $\pi$-accessible dans $M$, a fortiori accessible dans $M$, d'où a). Pour prouver que $\underline S'=\underline S$, par définition de $\underline S$ comme \unsure{la} $L$-enveloppe de $A$, il suffit de voir que $\underline S'$ est stable par \inserted{petites} $\underrightarrow{\lim}$ filtrantes. Or cela résulte du fait que $B$ \inserted{$\subset M_\pi$} est stable \inserted{dans $M_\pi$} par \deleted{\unsure{facteurs}} $\underrightarrow{\lim}$ filtrantes de cardinal $\leq\pi$, et de 4.13..19, qui implique tout au moins\,; \deleted{que} comme $B=\underline S'\cap M_\pi$ (en vertu de \unsure{4.13.4$_{\rm v}$|4.13.4$_{\rm iv}$}), que $\underline S'$ est stable par $\underrightarrow{\lim}$ filtrantes \uline{de cardinal} $\leq\pi$. Qu'il soit $=$ stable par \inserted{petites} $\underrightarrow{\lim}$ filtrantes quelconques, résulte de ceci, et du fait qu'il l'est par $\underrightarrow{\lim}$ grdes devant $\pi$. (cf 4.13.21\ldots)
% archive sheet number: 444
\scanpage{446}
\mspage{358}
\uline{Définition} 4.13.35. \unsure{Soient|Soit} $M$ une catégorie accessible, stable par petites $\underrightarrow{\lim}$ filtrantes, \deleted{Une partie} $S$ une partie de $\mathrm{Ob}\,M$, $\underline S$ la sous-catégorie \deleted{\unsure{str.}} pleine correspondante de $M$. On dit que $S$ est $L$-\uline{accessible}, \deleted{\unsure{si}} que $\underline S$ est $L$-accessible dans $M$, si $\underline S$ est accessible dans $M$, et de plus stable par \inserted{petites} $\underrightarrow{\lim}$ filtrantes.

\margin{\uline{NB} \deleted{Les} \inserted{L'ens des} parties \uline{$L$}-accessibles de $\mathrm{Ob}\,M$ est stable par petites intersections, et par réunions \uline{finies}. L'image inverse d'une partie \unsure{$L$|$L_0$}-accessible par un $L$-foncteur\textsuperscript{*} $f:M\unsure{\text{---}|\longrightarrow}N$ de cat. acc\,; \deleted{\unsure{et}} est $L$-accessible.}
% slanted
\margin{un \uline{$L_0$-foncteur} : qui commute aux petites $\underrightarrow{\lim}$ filtrantes.}

\uline{Proposition} 4.13.36. Soient $M,S,\underline S$ comme ds 4.13.35, et soit $\pi$ un cardinal utile pour $M$. \deleted{Conditions} \inserted{Supposons $\underline S$ \deleted{\unsure{$\pi$}} stable \deleted{par} dans $M$ par petites $\underrightarrow{\lim}$ filtrantes, ou} \deleted{équivalentes} seulement $\underline S(\pi)$ \inserted{$\overset{\mathrm{def}}{=}\underline S\cap M_\pi$} stable par petites $\underrightarrow{\lim}$ filtrantes \inserted{(\unsure{\deleted{dans}|dans} $M_\pi$)} de cardinal $\leq\pi$. Conditions équivalentes :

a) $\underline S$ est $\pi$-accessible, \deleted{et stable par petites $\underrightarrow{\lim}$ filtrantes} \deleted{\unsure{est est}} (\unsure{Il sera automatiquement|Il suivra toujours} stable par petites $\underrightarrow{\lim}$ filtrantes, i.e. $L$-accessible).

\deleted{a$'$) $\underline S$ est $\pi$-accessible, et $\underline S(\pi)\overset{\mathrm{def}}{=}\underline S\cap M_\pi$ est stable par $\underrightarrow{\lim}$ filtrantes de card $\leq\pi$, (propriété $L(\pi)$).}

b) \deleted{$\underline S(\pi)$ satisfait $L(\pi)$, et} Pour tt $x\in\mathrm{Ob}\,M$, on a $x\in\mathrm{Ob}\,\underline S$ sss $\underline S(\pi)_{/x}$ est cofinale ds $M_{\pi/x}$\unsure{,|.}

c) \deleted{Pour \unsure{tt} $\underline S(\pi)$ satisfait\,;} Pour tt $x\in\mathrm{Ob}\,M$, et tte représentation $x\simeq\underset{i\in I}{\underrightarrow{\lim}}\,x_i$
% archive sheet number: 445
\scanpage{447}
\mspage{359}
qui est $L$-$\pi$-adaptée (cf \unsure{page|pages} 343, c), on a ceci : $x\in\mathrm{Ob}\,S$ sss $I_{\underline S}$ est cofinal dans $I$.

c$'$) Comme c, avec $I$ \unsure{ordonnée|ordonné}.

\inserted{[ c$''$) (Si on a choisi, pour tt $x\in\mathrm{Ob}\,M$, \uline{une} représentation $\pi$-adaptée $x\simeq\underset{I}{\underrightarrow{\lim}}\,x_i$) : $x\in\mathrm{ob}\,\underline S\Longleftrightarrow$ \unsure{l'ens|alors} $I_{\underline S}$ est cofinal dans $I$ ]}

\margin{\deleted{\unsure{B}}}
\uline{Démonstration} C'est essentiellement l'énoncé de 4.13.22 dans le cas particulier où $\Phi={}$ ens des catégories \uline{filtrantes} de cardinal $\leq\pi$, mais où, pour simplifier, on a supposé d'avance que $\underline S(\pi)$ \unsure{satisfasse|satisfais|satisfait} à $L(\Phi)$. L'équivalence de a) b) c) est contenue dans \unsure{4.13.??|4.13.22}, ainsi que $c\Longrightarrow c'$. Mais \deleted{il faut} pour $c'\Longrightarrow a$, \deleted{\unsure{les|le}} l'hypothèse \inserted{accessoire} $\pi^{\unsure{\bullet|\circ|'}}=\pi$ et \deleted{le fait} que les $J\in\Phi$ sont grdes devant $\pi_0$, n'est pas satisfaite ici, et il faut une démonstration indépendante, calquée sur celle des \deleted{\unsure{les}} pages 333-336. \deleted{Dans \unsure{ce} résultat cette démonstration} Il suffit pour cel\unsure{a}, à l'endroit $\square$ page 333, d'invoquer le

\uline{Lemme} 4.13.37. \unsure{Soient|Soit} $M$ une \unsure{cat-|cat.}
% archive sheet number: 446
\scanpage{448}
\mspage{360}
accessible stable par $\underrightarrow{\lim}$ filtrantes, $\pi$ un cardinal utile pour $M$. Alors pour tt $x\in\mathrm{Ob}\,M$, on a une représentation
\[
x=\underset{i\unsure{\in|\notin} I}{\underrightarrow{\lim}}\,\unsure{x_i|x.}
\]
$L$-$\pi$-adaptée, [i.e. avec 1) $I$ équiv. à une petite catégorie et $I$ grde devant $\pi$, 2) $I$ stable par $\underrightarrow{\lim}$ filtrantes de cardin\unsure{al|aux} $\leq\pi$, 3) $I\unsure{\text{---}|\longrightarrow}M$ commute à ces $\underrightarrow{\lim}$, \unsure{4)|1)|d)} $x_i\in\mathrm{Ob}\,M_\pi\quad\forall i\in I$], avec de plus $I$ \uline{ordonnée}.

En effet, on part de
\[
x=\underset{j\in J}{\underrightarrow{\lim}}\,x_j,\qquad
J\cong(M_\pi/x)_0
\]
\inserted{\unsure{une} sous-catégorie réduite équiv. à $M_{\pi/x}$}

\deleted{\unsure{On pr}} On définit $I$ comme l'ens des sous-\unsure{catéories|catégories} \uline{filtrantes} \deleted{\unsure{simples}} \inserted{$J'$} de $J$, de \deleted{\unsure{cardinal}} \inserted{\deleted{\unsure{card}}} $\leq\pi$, \deleted{Il est imm} $I$ ordonné par l'inclusion \unsure{?|$\supseteq$|$\subseteq$}. \deleted{Il a} et on pose, pour $J'\in I$
\[
x_{J'}=\underset{j\in J'}{\underrightarrow{\lim}}\,x_j.
\]
Il est immédiat qu'on obtient ainsi un système inductif dans $M$, indexé par $I$.
% archive sheet number: 447
\scanpage{449}
\mspage{361}
D'autre part, il est clair que $I$ est petite, je dis que $I$ est stable par $\underrightarrow{\lim}$ filtrantes de cardinal $\leq\pi$, et que le foncteur $\deleted{x}J'\longmapsto x_{J'}:I\unsure{\text{---}|\longrightarrow}M$ y commute, enfin que les $x_{J'}$ sont dans $M_\pi$. Il reste à voir que $I$ est grde devant $\pi$, i.e. que \deleted{pour les sous-catégories} si on a $(J'_{\unsure{r|i}})_{r\in A}$, \inserted{avec} $\mathrm{Card}\,A\leq\pi$, les $J'_r$ dans $I$, alors $\exists\,J'\in I$ qui majore les $J'_r$. Soit \deleted{$K$} \inserted{$K$} \deleted{soit} la sous-catégorie \deleted{str. pleine} de $J$ engendrée par les $J'_r$, on voit \deleted{tout} \deleted{\unsure{facilement}} \deleted{qui est telle que $\mathrm{Ob}\,K=\bigcup\mathrm{Ob}\,J'_r$, \unsure{tout}} que $\mathrm{Card}\,K\leq\pi$. \unsure{restant|revenant} à prouver :

% slanted
\margin{\uline{NB} Pour la démonstration de 4.13.37, il suffit de prouver 4.13.38 dans le cas où $J$ est ordonné, auquel cas les \unsure{contorsions|complications|considérations} de la fin sont inutiles.}
\uline{Lemme} 4.13.38. Soient $J$ une catégorie gr\unsure{de|ande} devant $\pi$, $K$ une sous-catégorie de $J$ \deleted{avec $\mathrm{Card}$} de cardinal $\leq\pi$.

Alors il existe une sous-catégorie $K'$ \deleted{\unsure{cont}} \uline{filtrante} de $J$, contenant $K$ de cardinal $\leq\pi$, [\deleted{Si $K$ contient tous les objets \unsure{maxim$^{\rm x}$} de $J$,} \deleted{On peut $=$ préciser que} \deleted{$K'=K\cup\{e\}$ tel que $K'$ ait un objet final dans $K$} (Si on suppose que l'on \unsure{n'est|n'ait} pas, \deleted{on peut choisir \unsure{$K'$} tout entier \unsure{égal}} ds le cas où $J$ contient un objet maximal, et que $K$ contienne tous les objets maximaux de $J$],
% archive sheet number: 448
\scanpage{450}
\mspage{362}
\inserted{on peut $=$ \unsure{prendre|préciser}} $K'$ telle que $\mathrm{Ob}\,K'=\mathrm{Ob}\,K\unsure{\amalg|\sqcup|\cup}\{e\}$, où $e$ est un \uline{objet final} de $K'$, et $K$ \deleted{\uline{Dém}} \deleted{\unsure{étant}} \inserted{est} une sous-catégorie pleine de $K'$. (Donc \inserted{si $K'\ne K$,} $K'$ est isom. à un cône à droite construit sur $K$, ``déduit de $K$ en lui \unsure{ajoutant} un objet final''.)

\inserted{\uline{Dém} :} Si $K$ admet un objet final, on prend $K'=K$. Sinon, Comme $J$ est grde devant $\mathrm{Card}\,\mathrm{Ob}\,K$, il existe un $j\in\mathrm{Ob}\,J$ qui majore tous les objets de $K$. \deleted{\unsure{Choisissons pour}} \inserted{Choisissons}, \deleted{Pour tout $k\in\mathrm{Ob}\,K$, considérons alors flèches \unsure{toutes les}} pour tt $k\in\mathrm{Ob}\,K$, une flèche
\[
\alpha_k:k\longrightarrow j.
\]
Pour tte flèche
\[
u:k'\unsure{\text{---}|\longrightarrow}k
\]
de $K$, considérons la double-flèche
\[
(\alpha_k\circ u,\alpha_{\unsure{k'|m}}):k'\rightrightarrows j.
\]
L'ens de ces doubles flèches, de \deleted{\unsure{but}} $j$, est de \unsure{cardinal|cardinalité} $\leq\mathrm{Card}\,\mathrm{Fl}\,K\leq\pi$, donc $J$ étant grde devant $\pi$ il existe une flèche $\beta:j\unsure{\text{---}|\longrightarrow}e$ de $J$ qui égalise ttes ces doubles-flèches. Soit $K'\subset J$ ayant comme objets $\mathrm{Ob}\,K\cup\{e\}$, et comme \deleted{\unsure{couples}} flèches, \deleted{en plus} \unsure{En} plus de
% archive sheet number: 449
\scanpage{451}
\mspage{363}
celles de $K$, les flèches $\beta\alpha_k:k\longrightarrow e$ (pour $k\in\mathrm{Ob}\,K$), et $\mathrm{id}_e$. \uline{Si} $e\notin\mathrm{Ob}\,K$, on voit que c'est \inserted{bien} une sous-catégorie, admettant $e$ comme objet final et $K$\deleted{\unsure{*|s}} comme sous-catégorie pleine. Si on avait $e\in\mathrm{Ob}\,K$, mais il existe \deleted{\unsure{$e$ est dans}} \inserted{un objet $e'$} de $J$ qui majore $e$ et qui n'est pas dans $K$, \deleted{\unsure{par exemple|par contre} \unsure{on a|un} \unsure{Ob} \unsure{que}} on se ramène aisément au cas précédent, en remplaçant $e$ par $e'$ \inserted{\unsure{qui est un|qui est son} plus \unsure{gd|grand} \unsure{élément|élémt}}.
% slanted
\margin{La construction \unsure{qui|du} \deleted{\unsure{précède|point}} \unsure{suit|présent} est \unsure{déjà|en} contenue dans \unsure{l'énoncé|l'énonc.} plus général et plus précis 4.13.39}
\unsure{$[\![$|$[$} Donc il reste le cas où il existe \deleted{\unsure{un point|un objet} \unsure{de} $K$ \unsure{donc} \unsure{tel} \unsure{que} \unsure{tout objet} \unsure{de|objet}} \deleted{\unsure{tt objet de $K$},} \inserted{un objet $e$ dans $K$,} et \unsure{$e$|$e$'} tel que tout objet majorant $e$ soit dans $K$. \deleted{\unsure{Cela|tel}} Cela implique, $J$ étant filtrante, que $e$ est un \inserted{plus grand} objet \deleted{\unsure{maximal}} de $J$. (Car $\forall x\in J$, $\exists x'$ majorant commun de $x$ et $e$, donc $x'$ est dans $K$, donc majoré par $e$, qui majore donc $x$.) \unsure{=} On voit comme tantôt qu'il existe une flèche $\unsure{\beta|p}:e\longrightarrow j$ de $J$, qui égalise toutes les doubles flèches $k\rightrightarrows e$ de but $e$ \uline{dans $K$}. Quitte à composer avec $\unsure{j\longrightarrow e|j-e}$, ops $j=e$. On a donc un endom $p$ de $e$ dans $J$, tel que $pu=pv$ pour tt double-flèche \unsure{qq.|de|$\in$} \diagram{\begin{tikzcd}[ampersand replacement=\&,column sep=small] k\arrow[r,shift left,"u"]\arrow[r,shift right,"v"']\& e \end{tikzcd}} de $K$, de but $e$,
% archive sheet number: 450
\scanpage{452}
\mspage{364}
\deleted{\unsure{Alors|Donc}} \deleted{\unsure{Quel que soit|Ce qui revient à dire} \unsure{donc|pour} \unsure{que} \illegible{1}} Pour $k\in\mathrm{Ob}\,K$, soit
\[
 f_k:k\unsure{\longrightarrow|\text{---}}e
\]
la valeur commune des $pu$, pour $u\in\mathrm{Hom}_K(k,e)$. On voit tout de suite que si on a $k'\xrightarrow{v}k$ dans $K$, alors
\[
 f_{k'}=f_k\circ v.
\]
En particulier, prenant $k'=e$, on trouve
\[
 f_k\circ v=f_e\,(=p)\qquad\text{si }v:e\unsure{\text{---}|\longrightarrow}k.
\]
Soit donc \deleted{\unsure{$J$|$p$}} $K'$ la sous-catégorie de $J$ ayant \unsure{$=$|m\textsuperscript{es}} \unsure{objets|objet} que $K$, et comme flèches, en plus de celles de $K$, les \inserted{composées \diagram{\begin{tikzcd}[ampersand replacement=\&,column sep=small] k\arrow[r,"u"]\&e\arrow[r,"p^r"]\&e\arrow[r,"v"]\&k\unsure{''|'} \end{tikzcd}}} flèches \deleted{$f_k$ ($k$\unsure{ de|$\in$} $K$) $:k\longrightarrow e$. Il \unsure{faut}} \inserted{avec $u,v$ dans $K$, et $r\geq1$.} Il faut voir que c'est stable par composition (pour avoir bien une sous-catégorie).
% slanted
\margin{\unsure{Lire|faire} les flèches composées \diagram{\begin{tikzcd}[ampersand replacement=\&,column sep=tiny] k\arrow[r,"f_k"]\&e\arrow[r,"p^{r'}"]\&e\arrow[r,"v"]\&k' \end{tikzcd}}, avec $k$ et $k'\in\mathrm{Ob}\,K$, $v\in\mathrm{Fl}\,K$, $r'\geq0$.}
C'est clair pour \deleted{\diagram{\begin{tikzcd}[ampersand replacement=\&,column sep=small] k\arrow[r,"u"]\&k'\arrow[r,"{\unsure{v|u}}"]\&k'' \end{tikzcd}}} \unsure{ou|voir} si $u$ \deleted{\unsure{et|$v$}} $v$ sont dans $\mathrm{Fl}\,K$, \deleted{Si \unsure{$u$|$u$'} n'est pas ds $\mathrm{Fl}\,K$, c'est que l'on a $k'=e$, et $u=f_k$.} \unsure{Ou soit|En sorte|Il suffit}, on doit prendre un composé
\diagram{\begin{tikzcd}[ampersand replacement=\&,column sep=small]
 k\arrow[r,"u_1"]\&e\arrow[r,"p^r","{r\geq1}"']\&e\arrow[r,"u_2"]\arrow[rr,bend right=22,"{\unsure{p|\bullet}}"']\&k'\arrow[r,"v_1"]\&e\arrow[r,"p^s","{s\geq1}"']\&\deleted{\unsure{k'|k''|k}}e\arrow[r,"v_2"]\&k'',
\end{tikzcd}}
avec $u_1,u_2,v_1,v_2$ dans $\mathrm{Fl}\,K$, $r,s\geq1$. Or on a $p(v_1u_2)=p$, d'où
\[
 \begin{aligned}
 p^{\unsure{8|s|\deleted{r}s}}(v_1u_2)&=p^{s-1}(p\unsure{{}_1|{,}|.}v_1u_2)=p^{s-1}p\\
 &=p^s,
 \end{aligned}
\]
donc le composé est
\diagram{\begin{tikzcd}[ampersand replacement=\&]
 k\arrow[r,"u_1"]\&e\arrow[r,"p^{r+s}"]\&e\arrow[r,"v_2"]\&k'',
\end{tikzcd}}
que
% archive sheet number: 451
